Resume

(TR)

Bu bitirme projesinin amacı, farklı dillerdeki çevrim içi haberleri otomatik olarak toplayan, haberlerin ana metnini çıkaran, özetleyen ve kullanıcıya filtrelenebilir bir haber akışı olarak sunan yapay zekâ destekli bir web sistemi geliştirmektir. Proje, çevrim içi haber hacminin artmasıyla oluşan bilgi yoğunluğu problemini ele almaktadır. Kullanıcının dil, konu, bölge, ülke, kaynak ve anahtar kelime seçimlerine göre kısa ve doğrulanabilir özetlere hızlı biçimde ulaşmasını sağlar.

Projenin özgün değeri; haber toplama, metin çıkarımı, kalite kontrolü, çok dilli özetleme, isteğe bağlı çeviri, veritabanı tabanlı saklama ve web arayüzünü tek bir sistemde birleştirmesidir. Haberler RSS kaynaklarından toplanmış, metinler temizlenmiş ve XL-Sum veri setiyle eğitilen veya ince ayarlanan BART, mBART ve mT5 tabanlı yerel modellerle özetlenmiştir. Sistem ayrıca arka plan zamanlayıcısı, SQLite veritabanı, önbellekleme mekanizması ve kaynaklar arasında dengeli işleme yapısı kullanmaktadır. Modeller ROUGE, BLEU, METEOR-lite, sıkıştırma, gecikme, tutarlılık, doğruluk, açıklık, alaka ve verimlilik gibi ölçütlerle değerlendirilmiştir. Sonuç olarak 15 dili destekleyen ve özetlenmiş haberleri web arayüzü ile API üzerinden sunan çalışan bir prototip elde edilmiştir.

Gerçekçi kısıtlar tasarım sürecinde dikkate alınmıştır. Ekonomik kısıtlar, ücretli dış servislere bağımlılığı azaltan yerel modellerle sınırlandırılmıştır. Sürdürülebilirlik kısıtları, önbellekleme ve tekrarlı işlemlerin azaltılmasıyla ele alınmıştır. Etik, sosyal ve politik boyutlar için özetler özgün haber bağlantılarıyla birlikte sunulmuş, kaynak seçimi kullanıcıya bırakılmış ve kalite kontrolü uygulanmıştır. Proje; otomatik özetleme, çok dilli doğal dil işleme ve bilgiye hızlı erişim alanlarına katkı sağlayabilir. Ayrıca haber takibini kolaylaştırarak ve kullanıcıya zaman kazandırarak sosyo-ekonomik değer sunabilir.

---

(FR)

L’objectif de ce projet de fin d’études est de développer un système web assisté par intelligence artificielle qui collecte automatiquement des actualités en plusieurs langues, extrait le texte principal des articles, les résume et les présente sous forme d’un flux filtrable. Le projet répond au problème de la surcharge d’information créée par le grand volume d’actualités en ligne. Il permet à l’utilisateur d’accéder rapidement à des résumés courts et vérifiables selon la langue, le sujet, la région, le pays, la source et les mots-clés choisis.

La valeur originale du projet est de réunir dans un même système la collecte d’actualités, l’extraction de texte, le contrôle de qualité, le résumé multilingue, la traduction optionnelle, le stockage en base de données et une interface web. Les articles sont collectés à partir de sources RSS, puis les textes sont nettoyés et résumés avec des modèles locaux fondés sur BART, mBART et mT5, entraînés ou affinés avec le jeu de données XL-Sum. Le système utilise aussi un planificateur en arrière-plan, une base SQLite, un mécanisme de cache et un traitement équilibré des sources. Les modèles ont été évalués avec des mesures telles que ROUGE, BLEU, METEOR-lite, la compression, la latence, la cohérence, l’exactitude, la clarté, la pertinence et l’efficacité. Le résultat obtenu est un prototype fonctionnel qui prend en charge 15 langues et présente les actualités résumées via une interface web et une API.

Les contraintes réalistes ont été prises en compte dans la conception. Les contraintes économiques sont réduites par l’utilisation de modèles locaux, ce qui limite la dépendance aux services payants. Les contraintes de durabilité sont abordées par le cache et par la réduction des traitements répétés. Les aspects éthiques, sociaux et politiques sont pris en compte en affichant les liens vers les articles originaux, en laissant le choix des sources à l’utilisateur et en appliquant un contrôle de qualité. Ce projet peut contribuer aux domaines du résumé automatique, du traitement multilingue des langues et de l’accès rapide à l’information. Il peut aussi apporter une valeur socio-économique en facilitant le suivi de l’actualité et en faisant gagner du temps aux utilisateurs.